\documentclass[11pt]{article}
\usepackage{fontspec}
\usepackage[french]{babel}
\usepackage[leqno]{amsmath}
\usepackage{amssymb}
\usepackage{amsthm}
\usepackage{mathrsfs}
\usepackage{tikz-cd}
\usepackage{enumitem}

\newcommand{\printedpage}[1]{\marginpar{\footnotesize p.~#1}}
\newcommand{\msp}[1]{\marginpar{\footnotesize [#1]}}
\newcommand{\deq}{\overset{\text{déf}}{=}}
\newcommand{\Ho}{\mathrm{Ho}}
\newcommand{\Mono}{\mathrm{Mono}}
\newcommand{\Fl}{\mathrm{Fl}}
\newcommand{\Filt}{\mathrm{Filt}}
\newcommand{\Cof}{\mathrm{Cof}}
% printed pairing arrow (left harpoon down / right harpoon up on one stem)
\newcommand{\qpair}{\mathrel{\leftharpoondown\joinrel\relbar\joinrel\rightharpoonup}}

% footnote marks printed as superscript (n) in the text, plain superscript n in the footnote
\makeatletter
\renewcommand\@makefnmark{\hbox{\@textsuperscript{\normalfont(\@thefnmark)}}}
\renewcommand\@makefntext[1]{\parindent 1em\noindent\hbox{\@textsuperscript{\normalfont\@thefnmark}}#1}
\makeatother

\begin{document}

\printedpage{14}
% head: A. Grothendieck : Dérivateurs, Ch. XIX, Édition des Universités de Montpellier II et Paris VII
% foot: Version 11 février 2019

\msp{page 23}
L'axiome de décomposition est satisfait pour tout prédérivateur de la forme $\mathbf{D}_{\mathcal{M},W}$. L'axiome d'exactitude à gauche Dér 5 bis est valable, pourvu que les $(\mathcal{M}(I),W(I))$ s'insèrent dans une structure à cofibrations, ce qui est le cas. De même pour l'axiome Dér 2 de localisation, pourvu que de plus $W$ soit fortement saturé, ce qui est également le cas. Enfin, l'axiome de l'objet final dit que si $I$ admet un objet final, alors
\[ \Ho_W(e) \longrightarrow \Ho_W(I) \]
est pleinement fidèle, ou encore que le morphisme d'adjonction
\[ u_!u^*(\xi) \longrightarrow \xi \]
est un isomorphisme. Mais en vertu de 1.12, cas 3°, dans ce cas $u^*$ applique $\mathcal{M}(e)_{\mathrm{c}} = \mathcal{M}_{\mathrm{c}}$ dans $\mathcal{M}(I)_{\mathrm{c}}$, donc le calcul du couple de foncteurs adjoints $u_!^W$, $u_W^*$ peut se faire en utilisant les restrictions de $u_!$, $u^*$ à $\mathcal{M}(I)_{\mathrm{c}}$ et $\mathcal{M}_{\mathrm{c}}$ respectivement. Comme le morphisme d'adjonction $u_!^{\mathrm{c}}u_{\mathrm{c}}^* \longrightarrow \mathrm{id}$ est un isomorphisme, il en est donc de même de $u_!^W u_W^* \longrightarrow \mathrm{id}$, qui s'en déduit en localisant, q.e.d.

\textbf{2.5.} Pour achever de construire $\Ho_W$ comme un dérivateur sur Cat tout entier, il faut \msp{page 24} donc établir Dér 4, Dér 5 : existence des $u_*^W$, et calcul des fibres de $u_*^W$. Quant à Dér 6, il est acquis automatiquement, du fait que chaque $(\mathcal{M}(I),W(I))$ s'insère dans une structure de fibration $(W(I),F(I))$. (Oubli dans la démonstration de 2.4.) L'idée naturelle est d'utiliser à présent la structure de triple
\[ (W(I),C(I),\ \underbrace{F[I]}_{\deq TC(I)_*}\ )\ , \tag*{(2.5.1)} \]
mais la difficulté est de prouver la factorisation pour un couple
\[ C(I) \longleftrightarrow TF[I]\ , \qquad \underbrace{TC(I)}_{\deq C(I)\cap W(I)} \quad \longleftrightarrow F[I]\ , \tag*{(2.5.2)} \]
faute de pouvoir montrer, dans les conditions générales de 2.3 seulement (triple de Quillen spécial) que $C(I)$ et $TC(I)$ sont également spéciaux. (Cf. question 1.11.) Mais moyennant des conditions convenables sur $\mathcal{M}$, on arrivera à s'en tirer si $C = \mathit{mono}$. L'énoncé pertinent, un peu technique, se trouve dans XV cor. 7.4 (p. 98-99), que je vais rappeler ici.

\printedpage{15}
% head: A. Grothendieck : Dérivateurs, Ch. XIX, Édition des Universités de Montpellier II et Paris VII
% foot: Version 11 février 2019

\msp{page 25}
\textbf{Théorème 2.6.} \textit{Soit $\mathcal{M}$ une catégorie satisfaisant les conditions suivantes.}
\[ \left\{ \begin{minipage}{0.72\textwidth}\itshape
\begin{itemize}[leftmargin=2.5em]
\item[\upshape a)] $\mathcal{M}$ est “spéciale”, i.e. stable par petites $\varinjlim$ , et ses éléments sont accessibles.
\item[\upshape b)] Les limites inductives filtrantes dans $\mathcal{M}$ sont exactes à gauche (et dans $\mathcal{M}$ les $\varprojlim$ finies sont représentables – résulte de \textup{a} et \textup{c}).
\item[\upshape c)] $\mathcal{M}$ admet une filtration cardinale $(\Filt^\pi)_{\pi\geq\pi_0}$.
\item[\upshape d)] Les $\Filt^\pi$ sont stables par $\varprojlim$ indexées par des catégories $I$ telles que $\operatorname{card}\Fl(I) \leq \pi$, et par quotients ($Z \in \Filt^\pi$ et $Z \longrightarrow Z'$ épimorphisme implique $Z' \in \Filt^\pi$).
\item[\upshape e)] Si $A_1$, $A_2$ sont deux sous-objets d'un objet $B$ de $\mathcal{M}$, et $A_0 = A_1 \cap A_2$, alors $A_1 \amalg_{A_0} A_2 \longrightarrow B$ est un monomorphisme.
\item[\upshape f)] Si on a un carré cocartésien
\begin{center}
\begin{tikzcd}
A_0 \arrow[r, "i"] \arrow[d] & A_1 \arrow[d] \\
A_2 \arrow[r, "i'"] & A_3
\end{tikzcd}
\end{center}
avec $i$ monomorphisme, alors $i'$ est un monomorphisme.
\end{itemize}
\end{minipage}\right. \tag*{(2.6.1)} \]
$^{(10,11,12,13)}$\footnotetext[10]{\textbf{N.B.} a + c signifie aussi que $\mathcal{M}$ est une “catégorie $\mathfrak{U}$-paratopos”, et implique l'existence d'une filtration cardinale satisfaisant d) (sauf l'assertion concernant les quotients $Z'$ de $Z$). Donc il s'agit d'un $\mathfrak{U}$-paratopos satisfaisant de plus b) e) f). (Est-ce alors un topos\,? \textit{Non,} car e), f) sont O.K. dans les \textit{catégories abéliennes\,!})}\footnotetext[11]{\textbf{N.B.} La condition d) est sans doute inutile, du moins si les épimorphismes sont effectifs, car d) sera satisfaite en tous cas pour $\pi = 2^c$, avec $c \geq \pi_0$.}\footnotetext[12]{\textbf{N.B.} La condition e) n'est \textit{pas} satisfaite pour la catégorie (Gr), et le théorème est en défaut pour cette catégorie, laquelle n'admet \textit{pas} “assez d'injectifs”, cf. SGA 4 I 8.12.9. Par contre, e) et f) sont satisfaites par les catégories \textit{abéliennes.}}\footnotetext[13]{Cf. XVIII 4.6.7 (p. 157) pour une formulation plus sympathique de l'ensemble des conditions a) à f) pour $\mathcal{M}$.}. \textit{(On dira alors que $\mathcal{M}$ est} “très spéciale”\textit{). Soit de plus}
\[ W \ \subset \Fl(\mathcal{M}) \tag*{(2.6.2)} \]
\textit{un localiseur dans $\mathcal{M}$, qui satisfait les conditions}
\[ \left\{ \begin{minipage}{0.72\textwidth}\itshape
\begin{itemize}[leftmargin=2.5em]
\item[\upshape a)] $W$ est stable par $\varinjlim$ filtrantes $^{\textup{(14)}}$.
\item[\upshape b)] $W$ est de plus $\mathrm{L}_1$-accessible, i.e. il existe une \textup{petite} partie $W_0 \subset W$ telle que $W$ soit la $\mathrm{L}_1$-enveloppe de $W$, i.e. que $W$ soit la plus petite partie de $\Fl(\mathcal{M})$ contenant $W_0$, et stable par $\varinjlim$ filtrantes. % [sic] la L_1-enveloppe de W -> la L_1-enveloppe de W_0, W_0 being the small generating part just introduced
\end{itemize}
\end{minipage}\right. \tag*{(2.6.3)} \]
\footnotetext[14]{\textbf{N.B.} a) implique déjà que $W = W^\natural$.}

\printedpage{16}
% head: A. Grothendieck : Dérivateurs, Ch. XIX, Édition des Universités de Montpellier II et Paris VII
% foot: Version 11 février 2019

\msp{page 26}
\textit{Posons alors}
\[ C \ = \ \Mono(\mathcal{M}), \qquad TC \ = \ C\cap W\ , \qquad TF \ = \ C_*\ , \qquad F = (TC)_*\ . \tag*{(2.6.4)} \]
\textit{Sous ces conditions :}
\begin{itemize}[leftmargin=2.5em]
\item[1°)] \textit{$C \longleftrightarrow TF$, plus précisément $C \qpair TF$ est un couple de Quillen spécial.} % ?? uncertain: shape of the two arrows -- first: double-headed arrow (\longleftrightarrow) or harpoon pair; second: harpoon pair (left harpoon down / right harpoon up)
\item[2°)] \textit{Pour que $TC \qpair F$ soit un couple de Quillen, il faut et il suffit que $TC$ soit stable par cochangement de base, et alors c'est un couple de Quillen spécial.} % ?? uncertain: shape of the arrow -- harpoon pair (left harpoon down / right harpoon up)
\item[3°)] \textit{Pour que $(W,C,F)$ soit un triple de Quillen (lequel sera nécessairement spécial par \textup{1°, 2°}, et en particulier clos), il faut et il suffit que les conditions suivantes soient satisfaites.}
\[ \left\{ \begin{array}{ll}
\text{a)} & TC\ (\deq W\cap\Mono)\ \text{\itshape est stable par cochangement de base.}\\[1ex]
\text{b)} & TF\ (\deq (\Mono)_*) \ = \ \text{\itshape “applications de type injectif”}\ \subset W.
\end{array}\right. \tag*{(2.6.5)~$^{(15)}$} \]
\footnotetext[15]{\textbf{N.B.} Ces conditions (2.6.5) équivalent à dire que $(W,C)$ est une structure de catégorie à cofibrations sur $\mathcal{M}$.}
\textit{Et la condition} b) \textit{équivaut encore à :}
\[ \begin{minipage}{0.7\textwidth}
\begin{itemize}[leftmargin=2.5em]
\item[b$'$)] \textit{Toute flèche $f$ de $\mathcal{M}$ (ou seulement de $TF$) se factorise en $f = pi$, avec $i$ monomorphisme, $p \in W$.}
\end{itemize}
\end{minipage} \tag*{(2.6.6)} \]
\end{itemize}

\textbf{Corollaire 2.7.} \textit{$\mathcal{M}$ et $W$ étant comme dans} (2.6.1), (2.6.3), \textit{pour que $(W,C,F)$ soit un triple de Quillen} propre à gauche \textit{(i.e. $C \subset \Cof_W$), il faut et il suffit que $W$ satisfasse la condition} (2.6.5. b) \textit{(i.e. $W$ contient les flèches de type injectif), et de plus}
\msp{page 27}

\begin{itemize}[leftmargin=7em]
\item[] $C \subset \Cof_W$, \textit{i.e. la condition :}
\item[a$'$)] \textit{Pour tout diagramme cocartésien}
\[ \begin{tikzcd}
A_0 \arrow[r, "i"] \arrow[d, "f"'] & A_1 \arrow[d, "f'"'] \\
A_2 \arrow[r, "i'"] & A_3
\end{tikzcd} \tag*{(2.7.7)} \]
\textit{avec $i$ monomorphisme, si $f \in W$, alors $f' \in W$. (Stabilité de $W$ par cochangement de base par monomorphismes.)}
\end{itemize}

\textit{Sous ces conditions, ce triple de Quillen est spécial, et a fortiori clos.}

Ce corollaire résulte de 2.6. 3°), puisque $C \subset \Cof_W$ implique $TC \deq C\cap W = C\cap W^{\mathrm{univ}}$ (car $\Cof_W \cap W \ = \ W^{\mathrm{univ}}$), donc que $TC$ est stable par cochangement de base, d'où la condition (2.6.5. a).

\printedpage{17}
% head: A. Grothendieck : Dérivateurs, Ch. XIX, Édition des Universités de Montpellier II et Paris VII
% foot: Version 11 février 2019

\textbf{2.8.} On dit qu'une \textit{structure de Quillen} $(W,C,F)$ sur $\mathcal{M}$ est \textit{très spéciale,} si elle est obtenue comme au théorème 2.1., i.e. si on a les conditions :
\[ \left\{ \begin{minipage}{0.72\textwidth}
\begin{itemize}[leftmargin=2.5em]
\item[a)] $\mathcal{M}$ est très spéciale, i.e. satisfait aux conditions (2.6.1).
\item[b)] $W$ est stable par petites $\varinjlim$ filtrantes, et est accessible (conditions (2.6.3)).
\item[c)] $C = \Mono$.
\end{itemize}
\end{minipage}\right. \tag*{(2.8.1)} \]
Donc pour qu'un localiseur $W$ dans $\mathcal{M}$ corresponde à une structure de Quillen très spéciale \msp{page 28} dans $\mathcal{M}$ (laquelle sera forcément unique), il faut et il suffit que $\mathcal{M}$ et $W$ satisfassent les conditions (2.8.1. a)\,b)), et que de plus on ait :
\[ \left\{ \begin{minipage}{0.65\textwidth}
\begin{itemize}[leftmargin=2.5em]
\item[a)] $TC \deq \Mono\cap W$ est stable par cochangement de base.
\item[b)] $TF \deq (\Mono)_* \subset W$, i.e. $W$ contient les flèches de type injectif dans $\mathcal{M}$.
\end{itemize}
\end{minipage}\right. \tag*{(2.8.2)} \]
\footnote[16]{On dit aussi que le localiseur $W$ est très spécial, s'il définit une structure de Quillen très spéciale.}.

On trouve maintenant le

\textbf{Théorème 2.9.} \textit{Soit} $(W, C = \Mono, F = (TC)_* = (W\cap\Mono)_*)$ \textit{un triple de Quillen très spécial dans une catégorie $\mathcal{M}$} (2.8). \textit{Alors :}
\begin{itemize}[leftmargin=2.5em]
\item[1°)] \textit{Pour toute $I$ dans} Cat, $(W(I),C(I),F[I])$ \textit{est un triple de Quillen très spécial dans $\mathcal{M}(I)$, et on a en particulier}
\[ C(I) \ = \ \Mono(\mathcal{M}(I))\ , \tag*{(2.9.1)} \]
\textit{et de plus}
\[ F[I]\cap W(I) \ = \ (TF)[I]\ . \tag*{(2.9.2)} \]
\item[2°)] \textit{Pour toute flèche $u : I \longrightarrow I'$ dans} Cat, \textit{d'où le couple de foncteurs adjoints}
\[ \begin{tikzcd}
\mathcal{M}(I) \arrow[r, shift right, "u_*"'] & \mathcal{M}(I') \arrow[l, shift right, "u^*"']
\end{tikzcd}\ , \tag*{(2.9.3)} \]
\textit{on a}
\[ \left\{ \begin{array}{rcl}
u^*(W(I')) & \subset & W(I)\\
u^*(C(I')) & \subset & C(I)\\
u^*(TC(I')) & \subset & TC(I)\\
u_*(F[I]) & \subset & F[I']\\
u_*(TF[I]) & \subset & TF[I']\ ,
\end{array}\right. \tag*{(2.9.4)} \]
\textit{enfin $u_*$ induit un foncteur}
\printedpage{18}
% head: A. Grothendieck : Dérivateurs, Ch. XIX, Édition des Universités de Montpellier II et Paris VII
% foot: Version 11 février 2019
\msp{page 29}
\[ u_*^{\mathrm{f}} \ :\ \mathcal{M}(I)_{\mathrm{f}} \longrightarrow \mathcal{M}(I')_{\mathrm{f}} \qquad\quad ^{(17)} \tag*{(2.9.5)} \]
\footnotetext[17]{sous-catégories pleines des objets fibrants.}
compatible aux localiseurs $W(I)_{\mathrm{f}}$ \textit{et} $W(I')_{\mathrm{f}}$ \textit{induits par $W(I)$, $W(I')$ respectivement, donc induit un foncteur}
\[ u_*^W \ :\ \Ho_W(I) \longrightarrow \mathit{Ho}_W(I')\ . \tag*{(2.9.6)} \] % [sic] italic Ho_W(I') -> upright Ho_W(I'), inconsistent with the upright Ho_W elsewhere
\textit{Ce foncteur est un} adjoint à droite \textit{du foncteur}
\[ u_W^* \ :\ \mathit{Ho}_W(I') \longrightarrow \Ho_W(I) \tag*{(2.9.7)} \] % [sic] italic Ho_W(I') -> upright Ho_W(I'), inconsistent with the upright Ho_W elsewhere
\textit{induit par $u^*$} (2.9.3).
\end{itemize}
\textsc{Démonstration de 2.9.} Pour prouver 1°), il faut vérifier d'abord que les conditions préliminaires (2.8.1. a)\,b)) sur $\mathcal{M}$, $W$ sont stables par passage de $\mathcal{M}$, $W$ à $\mathcal{M}(I)$, $W(I)$, ce qui se fait par AQT \texttt{[âne qui trotte]}. D'ailleurs on vérifie tout de suite que
\[ C(I) \ = \ \Mono(\mathcal{M}(I))\ , \]
et le théorème 2.6 (1°, 2°) implique alors que
\[ \begin{array}{ccccc}
C(I) & \qpair & C(I)_* & \deq & TF[I]\\
\underbrace{TC(I)}_{\deq C(I)\cap W(I)} & \qpair & (TC(I))_* & \deq & F[I]
\end{array} \] % ?? uncertain: shape of the two arrows -- harpoon pair (left harpoon down / right harpoon up)
sont des couples de Quillen spéciaux. (En notant pour le deuxième, que $TC(I)$ est bien \msp{page 30} stable par cochangement de base, $TC$ l'étant.) Pour vérifier que l'on a un triple de Quillen très spécial, il reste en vertu de 2.8 à vérifier que
\[ TF[I] \ \subset\ W(I)\ , \]
ce qui s'écrit
\[ v^*(TF[I]) \ \subset\ W(I_0)\ , \]
où
\[ v \ :\ I_0 \longrightarrow I \]
est défini comme dans (1.12) (p. 10). En fait, on aura même
\[ v^*(TF[I]) \ = \ v^*v_*(TF(I_0)) \ \subset\ \underbrace{TF(I_0)}_{= F(I_0)\cap W(I_0)}\quad , \]
et la démonstration est duale de celle de
\[ v^*v_!(TC(I_0)) \ \subset\ TC(I_0) \]
(cf. (2.3.7), p. 21), elle ne dépend pas des questions de factorisation dans $\mathcal{M}(I)$, et provient ici de la stabilité de $TF$ par petits produits. Cela prouve 1°.

La démonstration de 2° est maintenant la duale de celle de (2.3. 2°). Comme cela nous fait répéter deux fois la même démonstration, cela montre qu'on n'a pas “monté le sorite” tout à fait avec assez de soin. Donc à titre de rattrapage, j'ajoute un

\end{document}
